\documentclass[]{article}
\usepackage{graphicx}
\usepackage[svgnames]{xcolor} 
\usepackage{fancyhdr}
\usepackage{fancyvrb}
\usepackage{forest}
\usepackage{tocloft}
\usepackage[hidelinks]{hyperref}
\usepackage{enumitem}
\usepackage[many]{tcolorbox}
\usepackage{listings }
\usepackage[a4paper, total={6in, 8in} , top = 2cm,bottom = 4cm]{geometry}
%\usepackage[a4paper, total={6in, 8in}]{geometry}
\usepackage{afterpage}
\usepackage{amssymb}
\usepackage{pdflscape}
\usepackage{textcomp}
\usepackage{xecolor}
\usepackage{rotating}
\usepackage[Kashida]{xepersian}
\usepackage[T1]{fontenc}
\usepackage{tikz}
\usepackage[utf8]{inputenc}
\usepackage{PTSerif} 
\usepackage{seqsplit}
\usepackage{changepage}


\usepackage{listings}
\usepackage{xcolor}
\usepackage{sectsty}

\setcounter{secnumdepth}{0}
 
\definecolor{codegreen}{rgb}{0,0.6,0}
\definecolor{codegray}{rgb}{0.5,0.5,0.5}
\definecolor{codepurple}{rgb}{0.58,0,0.82}
\definecolor{backcolour}{rgb}{0.95,0.95,0.92}
\definecolor{blanchedalmond}{rgb}{1.0, 0.92, 0.8}
\definecolor{brilliantlavender}{rgb}{0.96, 0.73, 1.0}
 
\NewDocumentCommand{\codeword}{v}{
\texttt{\textcolor{blue}{#1}}
}
\lstset{language=java,keywordstyle={\bfseries \color{blue}}}

\lstdefinestyle{mystyle}{
    backgroundcolor=\color{backcolour},   
    commentstyle=\color{codegreen},
    keywordstyle=\color{magenta},
    numberstyle=\tiny\color{codegray},
    stringstyle=\color{codepurple},
    basicstyle=\ttfamily\normalsize,
    breakatwhitespace=false,         
    breaklines=true,                 
    captionpos=b,                    
    keepspaces=true,                 
    numbers=left,                    
    numbersep=5pt,                  
    showspaces=false,                
    showstringspaces=false,
    showtabs=false,                  
    tabsize=2
}

\lstset{style=mystyle}

 \settextfont[BoldFont={XB Zar bold.ttf}]{XB Zar.ttf}


\setlatintextfont[Scale=1.0,
 BoldFont={LiberationSerif-Bold.ttf}, 
 ItalicFont={LiberationSerif-Italic.ttf}]{LiberationSerif-Regular.ttf}





\newcommand{\inputsample}[1]{
    ~\\
    \textbf{ورودی نمونه}
    ~\\
    \begin{tcolorbox}[breakable,boxrule=0pt]
        \begin{latin}
            \large{
                #1
            }
        \end{latin}
    \end{tcolorbox}
}

\newcommand{\outputsample}[1]{
    ~\\
    \textbf{خروجی نمونه}

    \begin{tcolorbox}[breakable,boxrule=0pt]
        \begin{latin}
            \large{
                #1
            }
        \end{latin}
    \end{tcolorbox}
}

\newtcolorbox{mybox}[2][]{colback=red!5!white,
colframe=red!75!black,fonttitle=\bfseries,
colbacktitle=red!85!black,enhanced,
attach boxed title to top center={yshift=-2mm},
title=#2,#1}

\newenvironment{changemargin}[2]{%
\begin{list}{}{%
\setlength{\topsep}{0pt}%
\setlength{\leftmargin}{#1}%
\setlength{\rightmargin}{#2}%
\setlength{\listparindent}{\parindent}%
\setlength{\itemindent}{\parindent}%
\setlength{\parsep}{\parskip}%
}%
\item[]}{\end{list}}


\definecolor{foldercolor}{RGB}{124,166,198}
\definecolor{sectionColor}{HTML}{ff5e0e}
\definecolor{subsectionColor}{HTML}{008575}

\definecolor{listColor}{HTML}{00d3b9}

\definecolor{umlrelcolor}{HTML}{3c78d8}

\definecolor{subsubsectionColor}{HTML}{3c78d8}

\defpersianfont\authorFont[Scale=0.9]{XB Zar bold.ttf}

\defpersianfont\titr[Scale=1.5]{Lalezar-Regular.ttf}

\defpersianfont\fehrest[Scale=1.2]{Lalezar-Regular.ttf}

\defpersianfont\fehrestTitle[Scale=3.0]{Lalezar-Regular.ttf}

\defpersianfont\fehrestContent[Scale=1.2]{XB Zar bold.ttf}


\sectionfont{\color{sectionColor}}  % sets colour of sections
\subsectionfont{\color{subsectionColor}}  % sets colour of sections
\subsubsectionfont{\color{subsubsectionColor}}


\renewcommand{\labelitemii}{$\circ$}


\renewcommand{\baselinestretch}{1.1}


\renewcommand{\contentsname}{فهرست}

\renewcommand{\cfttoctitlefont}{\fehrestTitle}


\renewcommand\cftsecfont{\color{sectionColor}\fehrestContent\selectfont}
\renewcommand\cftsubsecfont{\color{subsectionColor}\fehrestContent\selectfont}
\renewcommand\cftsubsubsecfont{\color{subsubsectionColor}\fehrestContent\selectfont}
%\renewcommand{\cftsecpagefont}{\color{sectionColor}}

\setlength{\parskip}{1.2pt}

\begin{document}


%%% title pages
\begin{titlepage}
\begin{center}

\textbf{ \Huge{به نام خدا} }
        
\vspace{0.2cm}

\includegraphics[width=0.4\textwidth]{sharif1.png}\\
\vspace{0.2cm}
\textbf{ \Huge{\emph درس برنامه‌سازی پیشرفته} }\\
\vspace{0.25cm}
\textbf{ \Large{ فاز سوم پروژه} }
\vspace{0.2cm}
       
 
      \large \textbf{دانشکده مهندسی کامپیوتر}\\\vspace{0.1cm}
    \large   دانشگاه صنعتی شریف\\\vspace{0.2cm}
       \large   ﻧﯿﻢ سال دوم 00-99 \\\vspace{0.10cm}
      \noindent\rule[1ex]{\linewidth}{1pt}
استاد:\\
    \textbf{{دکتر محمدامین فضلی}}



    \vspace{0.20cm}

   مهلت ارسال:\\
    \textbf{}
    \textbf{}

    \vspace{0.10cm}
مسئول پروژه:\\
    \textbf{\authorFont{امیرمهدی نامجو}}
    
        \vspace{0.10cm}
مسئول فاز دوم:\\
    \textbf{\authorFont{}}
    
        \vspace{0.10cm}
طراحان فاز دوم:\\
    \textbf{\authorFont{}}
    
        \vspace{0.05cm}
مسئولین تنظیم مستند:\\
    \textbf{\authorFont{پارسا محمدیان و سروش جهان‌زاد}}
    

\end{center}
\end{titlepage}
%%% title pages


%%% header of pages
\newpage
\pagestyle{fancy}
\fancyhf{}
\fancyfoot{}
\cfoot{\thepage}
\lhead{فاز سوم}
\rhead{\includegraphics[width=0.1\textwidth]{sharif.png}\\
دانشکده مهندسی کامپیوتر
}
\chead{پروژه برنامه‌سازی پیشرفته}
%%% header of pages
\renewcommand{\headrulewidth}{2pt}

\KashidaOff



\tableofcontents

\newpage

 \Large \textbf{\\\\
}


\section*{{\titr نکات قابل توجه}}
\addcontentsline{toc}{section}{{\fehrestContent نکات قابل توجه}}
\begin{itemize}
\item
پس از اتمام این فاز، در گیت خود یک تگ با ورژن \lr{"v3.0.0"} بزنید. در روز تحویل حضوری این tag بررسی خواهد شد و کدهای پس از آن نمره‌ای نخواهد گرفت. برای اطلاعات بیش‌تر در مورد شیوه ورژن‌گذاری، می‌توانید به
 \href{https://semver.org/}{\textcolor{blue}{\underline{این لینک}}}
 مراجعه کنید. البته برای این پروژه صرفا رعایت کردن همان ورژن گفته شده کافیست، اما خوب‌ است که با منطق ورژن‌بندی هم آشنا بشوید.

\item
در روز تحویل حضوری مشارکت تمام اعضای تیم در پروژه بررسی خواهد‌ شد و در صورت عدم مشارکت بعضی از اعضا، نمره‌ی ایشان برای آن فاز پروژه "صفر" لحاظ می‌گردد. مشارکت، با توجه به commit های افراد تیم در مخزن گیت‌هاب پروژه بررسی می‌شود.

\item
در هر فاز می‌توانید سه روز تاخیر به ازای کسر نمره داشته‌ باشید که به ازای هر روز آن، ۱۰ درصد از نمرهٔ آن فاز را از دست خواهید‌ داد. در مجموع سه‌فاز پروژه، سه روز تاخیر نیز بخشیده خواهد‌ شد.

\item
به ازای هر ساعتی که پروژه را زودتر تحویل دهید، ۱۵ دقیقه به مهلت تاخیر بدون کسر نمره شما اضافه خواهد‌ شد. این مقدار حداکثر یک روز خواهد‌ بود که در صورت ارسال ۴ روز زودتر از ددلاین به شما تعلق خواهد گرفت. \textbf{بنابراین ددلاین‌های پروژه تحت هیچ شرایطی تمدید نخواهد‌ شد.} توصیه می‌شود با برنامه‌ریزی مناسب به ددلاین‌های درس پایبند باشید.

\item
در صورت کشف تقلب از هریک از تیم‌ها، برای بار اول منفی نمرهٔ آن فاز برای آن تیم ثبت می‌شود و برای بار دوم، نمرهٔ منفی کل پروژه برای تیم لحاظ خواهد‌ شد که معادل مردود شدن در درس است.


\end{itemize}

\newpage

\section*{{\titr مقدمه}}
\addcontentsline{toc}{section}{{\fehrestContent مقدمه}}

همانطور که میدونید، در دو فاز قبلی بخش اصلی مورد نیاز برای اجرای بازی توسط یک کامپیوتر مستقل را پیاده‌سازی کردیم. یعنی بخش‌های مربوط به منطق و گرافیک بازی که برای کامل بودن بازی کافیه. پس در این بخش قراره چه کنیم؟ قراره قابلیت‌هایی رو به بازیمون اضافه کنیم که بشه بر بستر اینترنت هم بازی را اجرا و با بقیه بازیکنان از راه دور بازی کرد. همچنین این فاز بخش‌های امتیازی متنوع و زیادی داره که برای جبران نمرات از دست‌ رفته‌ی دو فاز قبلی می‌تونید ازشون استفاده کنید. در ادامه بخش‌های اصلی و امتیازی را به تفکیک می‌تونید ببینید.

\section*{{\titr بخش‌های اصلی}}
\addcontentsline{toc}{section}{{\fehrestContent بخش‌های اصلی}}

\subsection*{{\titr احراز هویت}}
\addcontentsline{toc}{subsection}{{\fehrestContent احراز هویت}}

فرض کنید می‌خواهید وارد گروهی شوید که برای افرادی که در آن ثبت‌نام می‌کنند قابلیت‌های ویژه‌ای در نظر گرفته شده است. شما ثبت‌نام کرده و مشخصات شما توسط مسئولین گروه ذخیره می‌شود، سپس به شما کارت عضویت داده می‌شود شما با این کارت احراز هویت شده و از شرایط خاص گروه بهره‌مند می‌شوید.

احراز هویت در برنامه‌نویسی نیز خاصیت مشابهی دارد شما پس از ثبت‌نام در سایتی دارای حساب کاربری شده و می‌توانید از ویژگی‌های خاص سایت مورد نظر استفاده کنید.

سوالی که مطرح می‌شود این است که احراز هویت به چه شکل صورت می‌گیرد؟

در ابتدا کاربر باید ثبت‌نام کند. به صورت عادی از کاربران نام‌کاربری و رمز عبور خواسته می‌شود (اطلاعات فرد)، کاربر پس از تایید ثبت‌نام دارای حساب کاربری شده و می‌تواند با نام‌کاربری و رمز عبور وارد حساب خود شود (کارت عضویت).

در برنامه‌‌ها سرور برای پردازش هر ریکوئست از سمت کلاینت نیاز به شناخت کاربر درخواست‌دهنده دارد.

عموما به دو شکل سرور متوجه کاربر حال حاضر می‌شود.

\begin{itemize}
    \item بر پایه‌ی سشن: \\
    در این روش سرور برای هر کاربری که لاگین می‌کند سشنی ساخته و رکوئست‌هایی که از آن دستگاه به سرور می‌آیند را در سشن مورد نظر پردازش می‌کند و متوجه کاربر درخواست دهنده می‌شود.
    \item بر پایه‌ی توکن: \\
    در این روش پس از لاگین‌کردن کاربر، سرور یک \lr{token} به کلاینت ارسال می‌کند. از این پس کلاینت با هر درخواست این توکن را به سمت سرور می‌فرستد تا سرور متوجه کاربر درخواست‌کننده شود.
\end{itemize}

پیشنهاد ما: استفاده از احراز هویت \lr{token-based}

از لینک‌های زیر می‌توانید برای مطالعه بیش‌تر استفاده کنید:

\begin{itemize}
    \item \href{https://dev.to/thecodearcher/what-really-is-the-difference-between-session-and-token-based-authentication-2o39}{تفاوت‌های این دو نوع}
    \item \href{https://sherryhsu.medium.com/session-vs-token-based-authentication-11a6c5ac45e4}{این دو در برابر هم}
    \item \href{https://www.okta.com/identity-101/what-is-token-based-authentication/}{احراز هویت بر پایه توکن}
\end{itemize}

\subsection*{{\titr اتاق گفت‌وگو (\lr{Chat Room})}}
\addcontentsline{toc}{subsection}{{\fehrestContent اتاق گفت‌وگو (\lr{Chat Room})}}

همانطور که در فاز یک اشاره شده بود، در این فاز قصد اضافه کردن اتاق گفت‌وگو به بازی خود را داریم.

هدف اصلی ما اضافه کردن اتاقی عمومی به برنامه به منظور صحبت و گفت‌وگو میان تمام افرادی است که در برنامه ثبت‌نام کرده‌اند.

\begin{tcolorbox}[colback=green!5!white,colframe=green!75!black,title=\textbf{نکته}]
    اضافه کردن «گفت‌و‌گوی خصوصی» و یا «گفت‌وگو در حین بازی» جزو خواسته‌های اصلی این قسمت نیست و هم‌چنین نمره امتیازی نیز نخواهد داشت؛ اما برای جذاب شدن هر چه بیشتر برنامه خود می‌توانید در صورت داشتن وقت و حوصله کافی :) این ویژگی‌ها را نیز به برنامه خود اضافه کنید.
\end{tcolorbox}

در بخش اصلی تنها باید
اتاق گفت‌وگوی عمومی(چت گلوبال \lr{Global Chat})
را پیاده‌سازی کنید.

\subsection*{{\titr تابلوی امتیازات (\lr{Scoreboard})}}
\addcontentsline{toc}{subsection}{{\fehrestContent تابلوی امتیازات (\lr{Scoreboard})}}

پیاده‌سازی تابلوی امتیازات باید به گونه‌ای باشد که اطلاعات آن در سمت سرور نگه داشته شده و هر بار کلاینت برای دیدن نتایج به سرور درخواست بدهد.

\subsection*{{\titr بازی دو نفره}}
\addcontentsline{toc}{subsection}{{\fehrestContent بازی دو نفره}}
در این بازی، بازیکنان باید بتوانند با یکدیگر به صورت‌های متفاوت (۱ راند و ۳ راند) با یکدیگر مسابقه دهند. به بیان دیگر، همان انجام بازی دو نفره که در فاز ۱ پیاده‌سازی کردید، در اینجا باید توسط کلاینت‌های متفاوت امکان پذیر باشد و هر بازیکن اتفاقاتی که می‌خواهد صورت بگیرد را در نوبت خودش به سمت سرور می‌فرستد و در صورت مجاز بودن خواسته‌ی بازیکن آن اتفاق عملی شود.

همچنین واضح است که بازیکنی که نوبتش نیست نیز باید به صفحه‌ی بازی دسترسی داشته باشد و به طور همزمان کارهایی که حریف انجام می‌دهد را ببیند. البته روشن است اگر بخش‌هایی از صفحه‌ی بازی لازم است برای حریف محدود باشد در اینجا نیز باید آن را رعایت کنید. مثلا کارت‌هایی که در دست حریف قرار دارند، نباید قابل مشاهده باشند.

\subsection*{{\titr لابی}}
\addcontentsline{toc}{subsection}{{\fehrestContent لابی}}
برای اینکه بازیکنان بتوانند همدیگر را برای مسابقه دادن با هم پیدا کنند باید یک لابی طراحی کنید در لابی بازیکنان می‌توانند با یکدیگر تعامل داشته باشند. (می‌توانید قسمت چت عمومی را نیز به لابی اضافه کنید) یکی از ویژگی‌های لابی اینست که بازیکن می‌تواند درخواست بازی دونفره با توجه به نوع بازی مد نظرش بدهد و در ادامه در حالت انتظار می‌ماند و اگر شخص دیگری هم درخواست مشابه با او را داشت سرور برای آنها یک بازی می‌سازد و آنها را به هم متصل می‌کند. توجه کنید، اگر چندین نفر هم زمان درخواست مشابه داشتند آن‌ها به صورت رندم به هم متصل می‌کند.

روشن است که بازیکنان باید بتوانند به اختیار خودشان از حالت انتظار خارج شوند.

\subsection*{{\titr فروشگاه}}
\addcontentsline{toc}{subsection}{{\fehrestContent فروشگاه}}
در این فاز فروشگاه بازی باید به سرور منتقل شود و تمامی کارت‌هایی که در شاپ وجود دارند از یک ظرفیت محدود برخوردار باشند. اگر بازیکنی کارتی را خرید، ظرفیت آن در شاپ یکی کم می‌شود. در صورتی که ظرفیت کارتی تمام شد هیچ بازیکنی، دیگر نمی‌تواند از آن کارت بخرد. تمامی تغییراتی که روی اقلام شاپ اتفاق می‌افتد باید در سرور ذخیره شود. از طرفی بازیکن‌ها می توانند کارت‌هایی که دارند را بفروشند و در این صورت ظرفیت آن کارت در شاپ یکی بیشتر می‌شود. همچنین برای فروشگاه باید یک پنل ادمین ساده برای مدیریت آن ایجاد شود. از طریق این پنل ادمین همچنین باید بتوان خرید و فروش یک کارت خاص را ممنوع کرد. به عنوان مثال اگر کارت \lr{X} به مقدار ۵تا در فروشگاه موجود باشد ولی از طریق پنل ادمین ممنوع شده باشد، اگر چه کارت در فروشگاه نمایش داده می‌شود اما باید یک علامت ممنوعه روی آن زده شده و امکان خریدن آن وجود نداشته باشد. از طریق این پنل کارهای ساده‌ی دیگر مانند افزودن یا کم کردن تعداد یک کارت باید قابل انجام باشد.




\section*{{\titr بخش‌های امتیازی}}
\addcontentsline{toc}{section}{{\fehrestContent بخش‌های امتیازی}}

\subsection*{{\titr اتاق گفت‌وگو (\lr{Chat Room})}}
\addcontentsline{toc}{subsection}{{\fehrestContent اتاق گفت‌وگو (\lr{Chat Room})}}

علاوه بر اتاق گفت‌و‌گوی عمومی، می‌توانید بخش‌های زیر را به صورت امتیازی پیاده‌سازی کنید.

\begin{itemize}
	\item قابلیت حذف پیام
	\item  قابلیت ویرایش (\lr{edit}) پیام
	\item قابلیت پاسخ دادن (\lr{reply}) پیام
	\item قابلیت سنجاق کردن (\lr{pin}) پیام
	\item نمایش آواتار ارسال کننده پیام
	\item نمایش اطلاعات کاربری فرد با کلیک بر روی آواتار او
	\item نمایش تعداد افراد آنلاین (آنلاین به این معنا که در حال حاضر در برنامه لاگین کرده و حضور دارند، نه این‌که لزوما در صفحه مربوط به اتاق گفت‌وگو باشند!).
\end{itemize}

\subsection*{{\titr دستاورد‌ها (\lr{Achievements})}}
\addcontentsline{toc}{subsection}{{\fehrestContent دستاورد‌ها (\lr{Achievements})}}

بازیکنان در طول اجرای بازی مهارت‌های مختلفی بدست خواهند آورد و تجربیات و کنش‌های مختلفی را پشت سر خواهند گذاشت. به پاس قدردانی از زحمات آن‌ها، نشان‌هایی برای آن‌ها در نظر بگیرید که اگر موفق به انجام کنش مورد نظر شدند، به آن‌ها تعلق گرفته و در پروفایل کاربری آن‌ها درج شود. این کنش‌ها را در ادامه مطرح می‌کنیم.

هم‌چنین علاوه بر نشان می‌توانید بسته به سختی این کنش‌ها مقداری پاداش برای هر عمل در نظر بگیرید که این پاداش می‌تواند اهدای یک کارت به بازیکن و یا اهدای مقداری پول و یا امتیاز به بازیکن باشد.

\textbf{\emph{تعریف کنش‌ها}}

در ادامه تعدادی کنش مطرح می‌شود که مقداردهی به پارامترهای آن‌ها پیشنهادی است اما به ذوق و سلیقه خودتان می‌توانید هر عدد معقولی بگذارید.

\begin{itemize}
	\item دنباله‌ای متوالی از پیروزی‌ها به تعداد مناسب‌ (مثلا ۵، ۱۰، ۱۵ و ...)
	\item پیروزی در تعداد مناسبی‌ (مثلا ۵) مسابقه فقط با استفاده از کارت‌های هیولا
	\item دنباله‌ای متوالی از پیروزی‌ها با محدودیت کاهش سلامتی مثلا سه برد متوالی با سلامتی بیش از ۳/۴
	\item پیروزی در تعداد محدودی نوبت
	\item دنباله‌ای متوالی از شکست‌ها به تعداد مناسب
	\item اولین - دهمین - صدمین و ... پیام اتاق گفت‌وگو
	\item احضار کارت به صورت \lr{Ritual Summon}
	\item ساخت اولین زنجیر
\end{itemize}


\begin{tcolorbox}[colback=green!5!white,colframe=green!75!black,title=\textbf{نکته}]
	دقت کنید که برخی از موارد بالا را می‌توان به صورت در لحظه \lr{(real time)} پیاده‌سازی کرد؛ به این معنا که بلافاصله بعد از برآورده شدن شرط مربوطه، دستاورد به کاربر تعلق گرفته و پیام مناسب نمایش داده شود؛ پیاده‌سازی این موارد به صورت ذکر شده نمره امتیازی خواهد داشت. امتیازدهی به کنش‌ها نیز متناسب با سختی پیاده‌سازی آن‌ها است.
\end{tcolorbox}

\subsection*{{\titr بخش \lr{TV}}}
\addcontentsline{toc}{subsection}{{\fehrestContent بخش \lr{TV}}}

صفحه‌ای را فرض کنید که دارای سه قسمت (تب) پخش زنده، پخش دوباره و پخش‌ برتر‌ها باشد.
از آنجا که منطق کلی هر بخش با بخش دیگر یکسان است (نمایش دادن یک بازی) می‌توان همه را در یک صفحه به نام صفحه TV گنجاند. در ادامه به شرح هر سه بخش می‌پردازیم.

\begin{itemize}
	\item \textbf{\emph{پخش بازی‌ها به صورت زنده}}
\end{itemize}
هدف این قسمت توانایی اشتراک‌گذاری آنلاین بازی‌هاست، هر کاربری که در حال بازی کردن می‌باشد باید بتواند بازی خود را به اشتراک بگذارد.
روش‌ پیشنهادی (ثبت \lr{Log}):

ابتدا باید حالت اولیه بازی ذخیره شود (کارت‌های هر بازیکن و اسپل‌ها و ...)

سپس هر حرکتی که صورت میگیرد باید ذخیره شود. می‌توان هر حرکت از بازی را ذخیره کرد، یعنی پس از هر حرکت جزئیات آن در قالب مشخصی ذخیره شود، سپس چون استیت اولیه بازی را می‌دانیم برای هر کاربری که قصد مشاهده را داشته باشد یک گیم ساخته و تمامی حرکات را براساس این گزارش‌ها را اجرا کرده تا به حالت فعلی در بازی برسیم.
این قسمت همانند استریم بازی‌ها است.

\begin{itemize}
	\item \textbf{\emph{مشاهده دوباره بازی}}
\end{itemize}
کاربر باید بتواند هر بازی خود را دوباره مشاهده و بررسی کند. 
به کمک روش ثبت \lr{Log} می‌توان هر بازی را ذخیره کرد و به حرکت قبلی و بعدی بازی دست یافت.


\begin{itemize}
	\item \textbf{\emph{مشاهده بازی‌های ضبط شده‌ی بازیکنان برتر}}
\end{itemize}
بازی بازیکنان برتر ذخیره شده و بتوان این بازی‌ها را مشاهده کرد.
به کمک ثبت \lr{Log} این موضوع حل می‌شود.

در نهایت می‌توانیم در کنار اسکوربرد بازی‌های اخیر بازیکنان برتر را نمایش دهیم.

\subsection*{{\titr امتیازدهی به بازی‌ها}}
\addcontentsline{toc}{subsection}{{\fehrestContent امتیازدهی به بازی‌ها}}

قابلیت امتیاز دادن به بازی‌های برتر، که در بخش تلویزیون برنامه نشان داده می‌شود، را به برنامه خود اضافه کنید.
\begin{itemize}
	\item داشتن سیستم امتیازدهی ساده
	\item به منظور قابل اطمینان بودن این امتیاز، باید به نوعی صلاحیت فرد امتیاز دهنده احراز شود مثلا حتما باید ۷۵ درصد از نوبت‌های بازی را تماشا کرده باشد یا هر آنچه که صلاح می‌دانید!
\end{itemize}

\subsection*{{\titr بازی دو نفره}}
\addcontentsline{toc}{subsection}{{\fehrestContent بازی دو نفره}}

\textbf{\emph{ساختن هوشمند بازی (\lr{Matchmaking}):}}

می‌توانید در منطق متصل کردن بازیکنان در سرور، منطقی قرار دهید که بازیکنانی با فاصله‌ی امتیازی قابل توجه به یکدیگر متصل نشوند یا با زمان انتظار بیشتری متصل شوند. با این حال باید برای زمان تحویل این بخش نیز سناریویی طراحی کنید.
به عبارتی تحویل گیرنده‌ی پروژه باید متوجه ۳ موضوع شود، در غیر این صورت امتیاز این بخش را به صورت کامل دریافت نمی‌کنید:

\begin{itemize}
	\item سرور بازیکنان با فاصله‌ی امتیازی را با تاخیر به هم متصل می‌کند.
	\item سرور بازیکنانی که امتیاز نزدیک به هم دارند را بدون وقفه به هم متصل می‌کند.
	\item اگر بازیکنان \lr{A} و \lr{B} با فاصله‌ی امتیازی قابل توجه در حال انتظار باشند و در این بین بازیکن C که امتیاز نزدیک به A دارد درخواست بازی داد، سرور A را به C متصل کند و B در حال انتظار بماند.
\end{itemize}

\subsection*{{\titr محدودیت زمانی هر نوبت}}
\addcontentsline{toc}{subsection}{{\fehrestContent محدودیت زمانی هر نوبت}}

هر نوبت (\lr{turn}) محدودیت زمانی خاصی دارد و اگر بازیکنی که نوبت اوست حرکت خود را انجام ندهد و زمان نوبت به پایان برسد بازی را می بازد و بازیکن رقیب برنده می‌شود.

\subsection*{{\titr امکان بازگشت به بازی بعد از قطع شدن اتصال کلاینت با سرور}}
\addcontentsline{toc}{subsection}{{\fehrestContent امکان بازگشت به بازی بعد از قطع شدن اتصال کلاینت با سرور}}

اگر در میانه ی یک turn ارتباط بازیکنی که نوبت اوست با سرور قطع شود تا زمان دوباره وصل شدن ارتباط بازیکن به سرور زمان turn ثابت نگه داشته می شود و بعد از اتصال دوباره زمان شروع به حرکت می‌کند.


\subsection*{{\titr مزایده در فروشگاه}}
\addcontentsline{toc}{subsection}{{\fehrestContent مزایده در فروشگاه}}

یکی دیگر از ویژگی‌هایی که باید در این فاز به بازی اضافه کنید مزایده است. مزایده به این صورت است که هر بازیکن می‌تواند از کارت‌هایی که در اختیار دارد یک کارت را با تعیین قیمت‌ اولیه‌ی آن برای فروش بگذارد. بازیکنان دیگر می‌توانند بیشترین هزینه‌ای را که تا به حال برای آن کارت پیشنهاد داده شده است (یا قیمت اولیه) را افزایش دهند در نهایت با اتمام زمانی که توسط سرور از قبل تعیین شده بازیکنی که بیشترین پیشنهاد را برای کارت داشته برنده‌ی مزایده می‌شود. در مورد نحوه‌ی کسر اعتبار از بازیکن برنده در ادامه توضیحاتی آورده‌ایم. برای همه‌ی مزایده‌ها مدت زمان مشخصی در نظر بگیرید همچنین هر بار که یک بازیکن قیمت پیشنهادی را افزایش میدهد زمان باقی‌مانده به مقدار مشخصی ریست شود. انتخاب این حداقل زمان و همچنین مدت زمان مزایده‌ها به سلیقه‌ی خودتان خواهد بود.

برای پیاده‌سازی این بخش علاوه بر اضافه کردن محلی برای شروع مزایده توسط هر بازیکن باید بخشی به بازی اضافه کنید که بازیکنان بتوانند مزایده‌های فعال را ببیند و وارد آن‌ها شود.

\textbf{\emph{ بخش اضافه کردن مزایده }}

در این بخش بازیکن به کارت‌های در اختیارش دسترسی دارد و می‌تواند یکی از آنها را انتخاب کند، سپس تنها با انتخاب یک قیمت اولیه مزایده را شروع می‌کند.

\textbf{\emph{ بخش مشاهده مزایده‌های فعال }}

باید در منو‌های هر بازیکن صفحه‌ای باشد که لیست تمامی مزایده‌های فعال همراه با آخرین قیمت پیشنهاد شده، زمان باقی مانده و عنوان کارت قابل مشاهده باشد، بازیکنان در این صفحه باید بتوانند پیشنهاد کارت‌ها را افزایش دهند. توجه کنید که با افزایش قیمت پیشنهادی کارت‌ها از اعتبار بازیکن به اندازه‌ی قیمتی که پیشنهاد داده کم می‌شود و اعتبار بازیکن قبلی که پیشنهاد بالا را داده بود به او برگردانده می‌شود. این صفحه باید یک دکمه برای رفرش کردن قیمت‌ها و زمان‌های باقی مانده داشته باشد. (در صورتی که بخش امتیازی را انجام نمی‌دهید)

\subsection*{{\titr نمایش بازیکنان آنلاین در اسکوربرد}}
\addcontentsline{toc}{subsection}{{\fehrestContent نمایش بازیکنان آنلاین در اسکوربرد}}

وضعیت آنلاین بودن بازیکنان باید در scoreboard مشخص شود. برای این‌کار باید یک لیست از سوکت‌هایی که بازیکنان با آن به سرور متصل هستند نگه‌داری کنید (احتمالا در بخش‌های دیگر پروژه نیز به این نیاز پیدا کرده‌اید/خواهید کرد) و هربار که شخصی درخواست مشاهده‌ی scoreboard را داشت سرور باید بازیکنانی که سوکت‌ آنها باز هست را آنلاین نشان دهد. همچنین در scoreboard یک دکمه‌ی رفرش قرار دهید که با استفاده از آن اطلاعات scoreboard از جمله وضعیت آنلاین بودن بازیکنان بروزرسانی شود.

\subsection*{{\titr امکان ارسال دعوتنامه برای بازی دو نفره}}
\addcontentsline{toc}{subsection}{{\fehrestContent امکان ارسال دعوتنامه برای بازی دو نفره}}

بازیکنان باید بتوانند برای یکدیگر دعوتنامه بفرستند با این توضیحات که یکی از ویژگی‌هایی که هر بازیکن می‌تواند در هنگام مشخص کردن نوع بازیش تعیین کند، اینست که می‌خواهد با چه کسی بازی کند. سپس مانند قبل منتظر پاسخ سرور می‌ماند. اگر بازیکنی که به او درخواست داده می‌شود این درخواست را بپذیرد سرور آن دو را به هم متصل می‌کند و اگر پیشنهاد را رد کند بازیکن درخواست دهنده از حالت انتظار خارج می‌شود. همچنین مانند قبل، بازیکن باید خودش نیز توانایی خارج شدن از حالت انتظار را داشته باشد در این حالت با توجه به سلیقه‌ی خودتان ممکن است که درخواست داده شده را حذف کنید یا هر امکان دیگری قرار دهید که اگر بازیکنی که به او درخواست داده می‌شود، درخواست را قبول کند متوجه منقضی شدن دعوتنامه شود.

\subsection*{{\titr نیاز نداشتن به رفرش}}
\addcontentsline{toc}{subsection}{{\fehrestContent نیاز نداشتن به رفرش}}
در بخش مشاهده‌ی مزایده‌های فعال و scoreboard اگر امکانی را فراهم کنید که بدون نیاز به دکمه‌ی رفرش کردن صفحه، اطلاعات صفحه بروزرسانی شود امتیاز این بخش را دریافت می‌کنید و همچنین ملزم به پیاده‌سازی دکمه‌ی رفرش نخواهید بود. به عبارتی:

\begin{itemize}
	\item در بخش مزایده‌های فعال: اگر بازیکن A در صفحه‌ی مزایده‌ها باشد و بازیکن B آیمبدون اینکه بازیکن A کاری انجام دهد باید مزایده‌ی اضافه شده را ببیند. همچنین اگر بازیکن C قیمت پیشنهادی برای مزایده‌ای را افزایش داد بازیکن A باید متوجه شود.
	\item در بخش scoreboard: اگر بازیکن A درصفحه‌ی scoreboard باشد و بازیکن B به بازی وارد شود بازیکن A باید متوجه این تغییر شود و همچنین اگر امتیاز بازیکن C افزایش یابد.
\end{itemize}

\subsection*{{\titr استفاده از دیتابیس}}
\addcontentsline{toc}{subsection}{{\fehrestContent استفاده از دیتابیس}}

پایگاه داده (database) مجموعه ای از داده‌ها است که به صورت منظم با ساختار خاصی درون کامپیوتر نگهداری شده است. پایگاه‌های داده معمولا توسط سامانه‌های مدیریت پایگاه داده (DBMS ها) کنترل می‌شوند. شما می‌توانید پایگاه داده خود را درون یک فایل به صورت ساده نگهداری کنید اما مهندسان بسیاری روی DBMS ها کار کرده‌اند و با توجه به شرایط کامپیوترهای موجود DBMS هایی را توسعه داده‌اند که عملیات های معمول مانند جستجو، دریافت داده‌هایی با شرایط خاص و تغییر در پایگاه داده را با سرعت بالا و راحتی انجام دهند. در ادامه با برخی از انواع DBMS ها و نحوه استفاده از آن ها آشنا می‌شویم.

\begin{itemize}
\item \textbf{\emph{ RDBMS }}
\end{itemize}

این خانواده از DBMS ها بسیار وسیع و در عین حال پر کاربرد هستند. RDBMS ها از زبان درخواست ساخت‌یافته (SQL) استفاده می‌کنند. در این دیتابیس‌ها داده‌ها در جدول‌هایی با ویژگی‌های مشخص ذخیره می‌شوند. بین این جداول روابطی (Relation) برقرار می‌شود و ما می‌توانیم عملیاتی مانند انتخاب (Select)، درج (Insert)، به روزرسانی (Update) و حذف (Delete) به همراه سایر ابزارهای مورد نیاز برای مدیریت پایگاه داده‌ها را در کنار هم داشته باشیم. RDBMS های معروف فراوانی وجود دارند و شما بسته به نیازتان از آنها استفاده می‌کنید. از جمله معروف‌ترین آن‌ها می توان به MySQL ،PostgreSQL ،SQLite و ... اشاره کرد.

\href{https://www.tutorialspoint.com/sqlite/sqlite_java.htm}{در این لینک میتوانید استفاده SQLite در جاوا را مشاهده کنید.}
\\

\begin{itemize}
\item \textbf{\emph{ NOSQL }}
\end{itemize}

با پیشرفت تکنولوژی و در نتیجه نیاز‌های بشر، دیگر RDBMS ها پاسخگوی همه‌ی نیازها نبودند. در نتیجه DBMS های NOSQL ساخته شد. این‌ها از زبان SQL پشتیبانی نمی کنند و در عوض قابلیت های دیگری مانند نگهداری پایگاه داده در چند سرور مجزا، جستجو در گراف و ... را دارند. Redis - Cassandra - MongdoDB - ElasticSearch  از دیتابیس های شناخته شده و پرکاربرد NOSQL هستند.

\href{https://www.tutorialspoint.com/mongodb/mongodb_java.htm}{در این لینک میتوانید استفاده MongoDB در جاوا را مشاهده کنید.}
\\

در این فاز شما محدودیتی در استفاده از دیتابیس‌ها ندارید و میتوانید از هر دیتابیس که می‌خواهید استفاده کنید.
\end{document}
